% -------------------------------------------------------------------------
% WBS ID: RES-VER-MET
% Deliverable: Test Metrics, Tolerances and Acceptance Criteria
% Status: In progress
% Evidence status: Local free-surface, conservation and loading metrics established
% Review status: Not reviewed
% Last updated: 2026-07-19
% Dependencies: RES-VER-FRM, RES-VER-BMK, RES-NUM-STB, RES-PHY-LOD
% Downstream interfaces: Software validation reports, RES-ML-DAT, RES-OPT-OBJ
% -------------------------------------------------------------------------

\section{RES-VER-MET --- Test Metrics, Tolerances and Acceptance Criteria}
\label{sec:res-ver-met}

\subsection{Purpose and Scope}
\label{sec:res-ver-met-purpose}

% TODO: Define what this Deliverable covers and explicitly excludes.

\subsection{Research Questions}
\label{sec:res-ver-met-research-questions}

% TODO: State the questions that must be answered before the Deliverable
% can be accepted.

\subsection{Terminology and Definitions}
\label{sec:res-ver-met-definitions}

% TODO: Define the technical terminology, symbols, quantities and
% classifications used in this Deliverable.

\subsection{Literature Evidence}
\label{sec:res-ver-met-literature-evidence}

% TODO: Record evidence from peer-reviewed papers, authoritative datasets,
% standards, technical reports and primary documentation.
%
% Do not create a detached literature-summary list. Explain how each source
% supports, contradicts or limits a project decision.

\textcite{qin2018comparison} supports validation against
depth-dependent velocity and direct force extraction. \textcite{watanabe2022validation}
supports multi-output acceptance because water level, velocity and
force can exhibit different error behaviour. Source role:
`3D-VALIDATION`. Project conclusion: no single scalar error metric is
sufficient for accepting the local URANS--VOF model.

\subsection{Governing Relations and Quantitative Definitions}
\label{sec:res-ver-met-governing-relations}

% TODO: Record relevant equations, dimensionless groups, empirical models,
% extraction rules or quantitative definitions.
%
% Use align environments for displayed equations.
% Every displayed equation must have a unique \label{eq:...}.
% Do not place punctuation after displayed equations.

\subsection{Test Objective}
\label{sec:res-ver-met-test-objective}

% TODO: Complete this RES-VER Domain-specific subsection.

\subsection{Reference or Benchmark Solution}
\label{sec:res-ver-met-reference-benchmark-solution}

% TODO: Complete this RES-VER Domain-specific subsection.

\subsection{Test Configuration}
\label{sec:res-ver-met-test-configuration}

% TODO: Complete this RES-VER Domain-specific subsection.

\subsection{Measured Quantities}
\label{sec:res-ver-met-measured-quantities}

% TODO: Complete this RES-VER Domain-specific subsection.

\subsubsection{Local-Impact Validation Observables}
\label{sec:res-ver-met-local-validation-observables}

The minimum validation set shall include:

\begin{itemize}
  \item phase-volume and mass-conservation diagnostics;
  \item $\alpha$ boundedness and interface-compression diagnostics;
  \item maximum local CFL and rejected-timestep count;
  \item pressure-correction residual and pressure-corrector count;
  \item free-surface elevation at multiple upstream and downstream
    locations;
  \item wave-front arrival time;
  \item depth-dependent velocity;
  \item inundation depth;
  \item overtopping volume or discharge;
  \item local pressure histories at several elevations;
  \item spatial pressure distribution;
  \item resultant force;
  \item resultant moment and force/moment history;
  \item force impulse;
  \item overturning moment;
  \item energy, reflection and transmission metrics;
  \item transmitted and reflected wave measures;
  \item wake and recirculation characteristics;
  \item structural displacement and restoring force where FSI is
    active.
\end{itemize}

The maximum value and complete time history shall both be assessed.
A model that reproduces one peak while misrepresenting load duration,
phase, or impulse shall not be accepted without further diagnosis.

\subsubsection{Benchmark Validation Quantities}
\label{sec:res-ver-met-benchmark-validation-quantities}

The benchmark shall report separately:

\begin{itemize}
  \item first-arrival time;
  \item leading-crest and leading-trough amplitude;
  \item waveform phase;
  \item dominant period and spectral content;
  \item maximum surface elevation;
  \item depth-integrated discharge or velocity;
  \item inundation depth;
  \item inundation extent;
  \item maximum run-up;
  \item reflected-wave behaviour;
  \item complete-wave-train duration.
\end{itemize}

A time-shifted waveform comparison may be used to assess waveform
shape, but the unshifted arrival-time error shall be reported first.
An imposed time shift shall not be treated as a correction to the
physical simulation.

\subsection{Error Metrics}
\label{sec:res-ver-met-error-metrics}

% TODO: Complete this RES-VER Domain-specific subsection.

\subsubsection{Free-Surface and Conservation Measures}
\label{sec:res-ver-met-free-surface-conservation-measures}

For free-surface elevation $\eta$, the normalised free-surface error is:

\begin{align}
  E_{\eta}
  &=
  \frac{
    \left\lVert
      \eta_{\mathrm{sim}}
      -
      \eta_{\mathrm{obs}}
    \right\rVert_{2}
  }{
    \left\lVert
      \eta_{\mathrm{obs}}
    \right\rVert_{2}
    +
    \epsilon_{\eta}
  }
  \label{eq:res-ver-met-free-surface-error}
\end{align}

The local mass-conservation acceptance metric is:

\begin{align}
  E_{\mathrm{mass}}
  &=
  \max
  \left(
    \max_i
    \widehat{R}_{m,i},\,
    \left|
      E_m
    \right|
  \right)
  \label{eq:res-ver-met-local-mass-acceptance}
\end{align}

where $\widehat{R}_{m,i}$ and $E_m$ are defined by the normalised
local mass residual and global mass-imbalance diagnostics in
`RES-NUM-STB` and `RES-NUM-ERR`.

VOF boundedness shall be reported as:

\begin{align}
  E_{\alpha,\mathrm{bound}}
  &=
  \max_i
  \left[
    \max
    \left(
      0,\,
      -\alpha_i,\,
      \alpha_i - 1
    \right)
  \right]
  \label{eq:res-ver-met-alpha-boundedness}
\end{align}

The maximum accepted-timestep Courant diagnostic shall be:

\begin{align}
  \mathrm{CFL}_{\max}
  &=
  \max_{f}
  \left(
    \frac{
      \left|
        \phi_f
      \right|
      \Delta t
    }{
      V_{c(f)}
      +
      \epsilon_V
    }
  \right)
  \label{eq:res-ver-met-cfl-maximum}
\end{align}

The pressure-correction residual shall be normalised as:

\begin{align}
  E_p
  &=
  \frac{
    \left\lVert
      \mathbf{A}_p p'
      -
      \mathbf{b}_p
    \right\rVert_2
  }{
    \left\lVert
      \mathbf{b}_p
    \right\rVert_2
    +
    \epsilon_p
  }
  \label{eq:res-ver-met-pressure-residual}
\end{align}

Decision role: these diagnostics are software-emitted acceptance
metrics for the RES-NUM-SPEC local solving procedure. Supporting
evidence: \textcite{watanabe2022validation} and \textcite{qin2018comparison}
support multi-output local validation; the numerical residual and CFL
definitions are project acceptance diagnostics tied to the selected
pressure-correction/VOF procedure. Limitation: final tolerance values
remain open.

\subsubsection{Local-Impact Error Measures}
\label{sec:res-ver-met-local-error-measures}

For a measured quantity $q_{\mathrm{obs}}(t)$ and simulated quantity
$q_{\mathrm{sim}}(t)$, the time-series root-mean-square error is:

\begin{align}
  \operatorname{RMSE}_{q}
  &=
  \sqrt{
    \frac{1}{N}
    \sum_{n=1}^{N}
    \left[
      q_{\mathrm{sim}}(t_n)
      -
      q_{\mathrm{obs}}(t_n)
    \right]^2
  }
  \label{eq:res-ver-met-local-rmse}
\end{align}

The peak-value ratio is:

\begin{align}
  R_{q,\max}
  &=
  \frac{
    \max_{t}
    \left|
    q_{\mathrm{sim}}(t)
    \right|
  }{
    \max_{t}
    \left|
    q_{\mathrm{obs}}(t)
    \right|
  }
  \label{eq:res-ver-met-local-peak-ratio}
\end{align}

The normalised impulse error is:

\begin{align}
  E_{\mathrm{J}}
  &=
  \frac{
    \left\lVert
    \mathbf{J}_{\mathrm{sim}}
    -
    \mathbf{J}_{\mathrm{obs}}
    \right\rVert
  }{
    \left\lVert
    \mathbf{J}_{\mathrm{obs}}
    \right\rVert
    +
    \epsilon_{\mathrm{J}}
  }
  \label{eq:res-ver-met-local-impulse-error}
\end{align}

The normalised moment error is:

\begin{align}
  E_{\mathrm{M}}
  &=
  \frac{
    \left\lVert
    \mathbf{M}_{\mathrm{sim}}
    -
    \mathbf{M}_{\mathrm{obs}}
    \right\rVert
  }{
    \left\lVert
    \mathbf{M}_{\mathrm{obs}}
    \right\rVert
    +
    \epsilon_{\mathrm{M}}
  }
  \label{eq:res-ver-met-local-moment-error}
\end{align}

The overtopping-volume error is:

\begin{align}
  E_{\mathrm{ot}}
  &=
  \frac{
    \left|
      V_{\mathrm{ot,sim}}
      -
      V_{\mathrm{ot,obs}}
    \right|
  }{
    \left|
      V_{\mathrm{ot,obs}}
    \right|
    +
    \epsilon_{\mathrm{ot}}
  }
  \label{eq:res-ver-met-overtopping-error}
\end{align}

The run-up error is:

\begin{align}
  E_{\mathrm{ru}}
  &=
  \frac{
    \left|
      R_{\mathrm{u,sim}}
      -
      R_{\mathrm{u,obs}}
    \right|
  }{
    \left|
      R_{\mathrm{u,obs}}
    \right|
    +
    \epsilon_{\mathrm{ru}}
  }
  \label{eq:res-ver-met-runup-error}
\end{align}

The energy-diagnostic error for any selected energy metric $E_\star$
is:

\begin{align}
  E_{\mathrm{energy},\star}
  &=
  \frac{
    \left|
      E_{\star,\mathrm{sim}}
      -
      E_{\star,\mathrm{ref}}
    \right|
  }{
    \left|
      E_{\star,\mathrm{ref}}
    \right|
    +
    \epsilon_{\star}
  }
  \label{eq:res-ver-met-energy-metric-error}
\end{align}

Source role: this metric provides a common acceptance form for the
incident, reflected, transmitted and apparent-dissipation measures in
RES-PHY-BAR. Supporting evidence: \textcite{watanabe2022validation}
supports separate metric reporting rather than one aggregate score.
Limitation: the reference energy measure may be experimental,
high-fidelity numerical or comparative depending on the V1--V2--V3
stage.

Force and moment extraction shall also pass surface-integration
consistency checks by comparing integrated patch traction with
component pressure and shear contributions:

\begin{align}
  E_{\mathrm{F,int}}
  &=
  \frac{
    \left\lVert
      \mathbf{F}
      -
      \mathbf{F}_{p}
      -
      \mathbf{F}_{\tau}
    \right\rVert
  }{
    \left\lVert
      \mathbf{F}
    \right\rVert
    +
    \epsilon_{\mathrm{F}}
  }
  \label{eq:res-ver-met-force-integration-consistency}
\end{align}

\subsection{Tolerance and Acceptance Criteria}
\label{sec:res-ver-met-tolerance-acceptance-criteria}

% TODO: Complete this RES-VER Domain-specific subsection.

\subsubsection{Local-Model Acceptance Principle}
\label{sec:res-ver-met-local-acceptance-principle}

The operational URANS model shall be accepted only where it satisfies
both:

\begin{enumerate}
  \item experimental validation requirements for the quantities
    relevant to barrier performance and structural loading;
  \item controlled agreement with the selected high-fidelity model
    for representative candidate geometries.
\end{enumerate}

Neither URANS nor LES shall be accepted solely from its theoretical
fidelity. The selected formulation must demonstrate that it preserves:

\begin{itemize}
  \item the relevant physical mechanisms;
  \item the required engineering outputs;
  \item the relative performance ranking of candidate barriers;
  \item conservative and reproducible load transfer to the structural
    model.
\end{itemize}

The final acceptance tolerances shall account for experimental
uncertainty, measurement filtering, repeatability, and the intended use
of each model-fidelity level.

Acceptance for the baseline local 3D model shall include:

\begin{itemize}
  \item free-surface elevation and arrival-time error below tolerance;
  \item VOF boundedness, phase volume and mass-conservation
    diagnostics below tolerance;
  \item CFL, pressure residual and rejected-timestep counts below
    tolerance;
  \item pressure, force-history, impulse and moment-history errors
    below tolerance;
  \item run-up and overtopping errors below tolerance;
  \item incident, reflected, transmitted and dissipated-energy metrics
    reported and below the stage-appropriate tolerance where reference
    data exist;
  \item no unresolved boundary-reflection contamination;
  \item stable candidate ranking under mesh, timestep and URANS--LES
    comparison studies.
\end{itemize}

\subsection{Execution Handoff}
\label{sec:res-ver-met-execution-handoff}

% TODO: Complete this RES-VER Domain-specific subsection.

\subsection{Candidate Approaches}
\label{sec:res-ver-met-candidate-approaches}

% TODO: Define the credible candidate approaches considered.

\subsection{Critical Comparison}
\label{sec:res-ver-met-critical-comparison}

% TODO: Compare candidates using physically and project-relevant criteria.
% Do not select an approach solely on popularity or implementation ease.

\subsection{Assumptions and Applicability}
\label{sec:res-ver-met-assumptions}

% TODO: State assumptions and the conditions under which they are valid.

\subsection{Limitations and Failure Conditions}
\label{sec:res-ver-met-limitations}

% TODO: State where the evidence, model or selected approach ceases to be
% reliable or applicable.

\subsection{Project-Specific Conclusions}
\label{sec:res-ver-met-conclusions}

% TODO: Convert the evidence into conclusions specific to the Studentship
% project. Do not merely restate literature findings.

The local 3D model is valid as a baseline high-Reynolds-number
hydrodynamic tsunami-barrier interaction model for run-up, overtopping,
free-surface deformation, pressure, force, moment, and comparative
barrier-performance analysis. It is not valid yet for structural
failure, scour, deformable barriers, detailed aeration, sediment
transport, or final design certification.

\subsection{Selected Approach and Rationale}
\label{sec:res-ver-met-selected-approach}

% TODO: Record the selected approach only after sufficient evidence exists.
% State the alternatives rejected and the reasons for rejection.

The selected acceptance approach is output-specific. Free-surface,
mass-conservation, pressure, force, impulse, moment, run-up and
overtopping metrics shall be reported separately rather than collapsed
into one aggregate validation score.

\subsection{Unresolved Decisions}
\label{sec:res-ver-met-unresolved-decisions}

% TODO: Record unresolved questions, required evidence and decision owners.

Open items are final numerical tolerance values, observation datasets,
experimental uncertainty models, acceptable URANS--LES discrepancy,
ranking-stability threshold and certification-level design criteria.

\subsection{Requirements and Handoffs}
\label{sec:res-ver-met-requirements}

% TODO: State requirements passed to other Research Domains, Software,
% verification activities or later execution work.

Software validation reports shall provide the metrics in this
Deliverable for each accepted local run, including mass conservation,
$\alpha$ boundedness, CFL, pressure residual, force/moment histories,
run-up, overtopping and energy metrics. `RES-ML-DAT` shall use only
validated outputs for any future training data, and `RES-OPT-OBJ`
shall not treat unvalidated hydrodynamic metrics as final design
objectives.

\subsection{Deliverable Completion Record}
\label{sec:res-ver-met-completion-record}

% TODO: Record whether the Deliverable is:
%
% - Not started
% - In progress
% - Draft complete
% - Reviewed
% - Accepted
%
% Acceptance requires:
%
% - the research questions to be answered;
% - claims to be supported by appropriate evidence;
% - assumptions and limitations to be explicit;
% - the project-specific conclusion to be stated;
% - downstream requirements to be traceable;
% - unresolved decisions to be closed or formally deferred.
